\documentclass[a4paper,11pt]{article}
\usepackage[T2A]{fontenc}
\usepackage[utf8]{inputenc}
\usepackage[russian]{babel}
\usepackage{amsmath,amssymb,mathtools}
\usepackage{geometry}
\geometry{left=20mm,right=20mm,top=20mm,bottom=21mm}
\usepackage{microtype}
\usepackage{array,tabularx,booktabs}
\usepackage{enumitem}
\usepackage{xcolor}
\usepackage{hyperref}
\usepackage{tikz}
\definecolor{accent}{HTML}{008080}
\definecolor{deepaccent}{HTML}{00575A}
\definecolor{ink}{HTML}{30343B}
\definecolor{err}{HTML}{B22B2B}
\hypersetup{colorlinks=true,linkcolor=deepaccent,urlcolor=deepaccent,pdfauthor={Лёвин Артём Александрович},pdftitle={Ревью домашней работы — Ксения Васильченко}}
\setlength{\parindent}{0pt}
\setlength{\parskip}{5pt}
\setlist[itemize]{topsep=3pt,itemsep=2pt,leftmargin=16pt}
\setlength{\emergencystretch}{2em}
\newcommand{\block}[1]{\par\vspace{5pt}{\bfseries\color{deepaccent}#1}\par\vspace{1pt}}
\newcommand{\problem}[2]{\clearpage\phantomsection\label{task:#1}{\Large\bfseries\color{deepaccent}Задание №#1}\hfill{\small #2}\par\vspace{3pt}\color{ink}}
\newcommand{\errorword}{\textcolor{err}{\textbf{Ошибка:}}}
\newcommand{\idea}{\textcolor{deepaccent}{\textbf{Ключевая идея:}}}
\newcommand{\outcome}{\textcolor{deepaccent}{\textbf{Итог:}}}
\begin{document}
\color{ink}
\thispagestyle{empty}
\begin{center}
\vspace*{29mm}
{\small\color{deepaccent}МАТЕМАТИКА}\par
\vspace{19mm}
{\LARGE Ревью домашней работы}\par
\vspace{16mm}
{\Large\color{accent}Доли, отношения и пропорции}\par
\vspace{12mm}
{\normalsize Разбор ошибок и самостоятельное закрепление}\par
\vfill
{\large Ксения Васильченко}\par
\vspace{8mm}
{18 ноября 2026 г.}\par
\vspace{17mm}
{\small Лёвин Артём Александрович}\par
{\small эксклюзивно для Ксении Васильченко}\par
\vspace{19mm}
\begin{tikzpicture}
\draw[accent,line width=0.8pt] (0,0) .. controls (2,0.65) and (3,-0.4) .. (4.5,0) .. controls (6,0.45) and (7,-0.65) .. (9,0);
\end{tikzpicture}
\end{center}
\clearpage
\section*{Сводная таблица ошибок}
\vspace{2mm}
\small
\renewcommand{\arraystretch}{1.36}
\begin{tabularx}{\linewidth}{@{}p{9mm}p{28mm}X p{27mm}p{27mm}@{}}
\toprule
№ & Тема & Суть недочёта & Ваш ответ & Верный ответ \\
\midrule
\hyperref[task:1]{1} & Доля от целого & Вычисление не доведено до ответа на вопрос & $\frac9{12}:\frac34$ & $\frac34$ \\
\hyperref[task:3]{3} & Доля в совокупности & В доле белых использовано число красных & $\frac35$ & $\frac25$ \\
\hyperref[task:5]{5} & Отношение масс & Не согласованы единицы; порядок величин обратный & $\frac1{25}$ & $\frac1{40}$ \\
\hyperref[task:7]{7} & Пропорция & Не записан формат $1:n$; пункт «б» не решён & $\frac17$; далее нет ответа & $1:7$; $16{,}8$ кг \\
\hyperref[task:8]{8} & Сравнение долей & Сравнены неудачи вместо успехов & Лена & Маша \\
\hyperref[task:9]{9} & Условная доля & Решение отсутствует & Не указан & $60\%$ \\
\hyperref[task:10]{10} & Сравнение долей & Число неудач принято за число успехов & A; $\frac54$ & Б; $\frac{16}{15}$ \\
\bottomrule
\end{tabularx}
\normalsize
\par\vspace{8mm}
\noindent\textbf{Как работать с ревью.} Откройте нужное задание по номеру в таблице. В каждом разборе найдите первый неверный или пропущенный шаг, сравните его с исправлением и выполните одно новое упражнение. Правильно выполненные задания не дублируются.
\clearpage
\thispagestyle{empty}
\vspace*{0.34\textheight}
\begin{center}
{\LARGE\bfseries\color{deepaccent}Разбор ошибок}
\end{center}

\problem{1}{Доля от общего количества}
\block{Текст задания}
Из 12 задач решены 9. Какую часть всех задач составляют решённые?
\block{Ваше решение и ответ}
Записано: $\frac9{12}:\frac34$. Заключительный ответ на вопрос не указан.
\block{Где именно возникает недочёт}
Последний корректный шаг --- выбор чисел 9 (решено) и 12 (всего). Далее появляется деление $\frac9{12}:\frac34$, хотя для нахождения доли оно не требуется. Это не даёт итоговой части всех задач.
\block{Почему это неверно}
\errorword\ Доля показывает, сколько единиц нужного вида приходится на одну единицу целого. Поэтому необходимо \emph{разделить количество решённых задач на количество всех задач}. Деление найденной доли на $\frac34$ отвечает уже на иной вопрос: «во сколько раз одна доля больше другой»; здесь оно равно 1, но 1 не является искомой долей.
\block{Ключевая идея}
\idea\ Часть целого вычисляется по формуле
\[
\text{доля}=\frac{\text{количество нужных объектов}}{\text{общее количество объектов}}.
\]
\block{Исправление и полное правильное решение}
Всего 12 задач; из них решены 9. Значит,
\[
\frac9{12}=\frac{9:3}{12:3}=\boxed{\frac34}.
\]
В числителе и знаменателе разделили на 3, поэтому значение дроби не изменилось.
\block{Проверка}
Три четверти от 12 равны $12\cdot\frac34=9$ задачам --- условие выполняется.
\block{Итог}
\outcome\ Решённые задачи составляют $\boxed{\frac34}$ всех задач.
\block{Задача на закрепление}
Из 20 упражнений выполнены 15. Какую часть всех упражнений составляют выполненные?

\problem{3}{Часть и отношение}
\block{Текст задания}
В букете 30 красных и 20 белых тюльпанов. Найдите: а) отношение числа красных к числу белых; б) долю белых среди всех тюльпанов.
\block{Ваше решение и ответ}
В пункте «а» записано $\frac{30}{20}=\frac32$ --- верно. В пункте «б» найдено $20+30=50$, затем использовано $\frac{30}{50}=\frac35$.
\block{Где именно возникает ошибка}
Последний верный шаг: $20+30=50$ --- это общее число цветов. Первый неверный шаг: $\frac{30}{50}$ как доля \emph{белых}. Число 30 обозначает красные тюльпаны, а не белые.
\block{Почему это неверно}
\errorword\ В знаменателе дроби должно стоять число всех тюльпанов, а в числителе --- именно тех, долю которых спрашивают. Выражение $\frac{30}{50}$ правильно вычисляет долю \emph{красных}, но отвечает на другой вопрос. Даже математически правильное деление не исправляет неверный выбор числителя.
\block{Ключевая идея}
\idea\ Сначала определите, какой объект назван в вопросе. Только после этого выбирайте числитель: «доля белых» $\Rightarrow$ в числителе 20.
\block{Исправление и полное правильное решение}
\textbf{а)} Отношение красных к белым:
\[
30:20=3:2 \quad\left(\frac{30}{20}=\frac32\right).
\]
\textbf{б)} Всего $30+20=50$ тюльпанов. Белых --- 20, следовательно,
\[
\frac{20}{50}=\frac{20:10}{50:10}=\boxed{\frac25}.
\]
\block{Проверка}
Доли красных и белых должны в сумме дать единицу:
\[
\frac35+\frac25=\frac55=1.
\]
\block{Итог}
\outcome\ а) $3:2$; б) $\boxed{\frac25}$ всех тюльпанов --- белые.
\block{Задача на закрепление}
В коробке 18 синих и 12 зелёных карандашей. Найдите долю зелёных карандашей среди всех карандашей.

\problem{5}{Отношение масс}
\block{Текст задания}
Масса котёнка --- 120 г, кошки --- 4,8 кг. Найдите отношение массы котёнка к массе кошки.
\block{Ваше решение и ответ}
Записано $\frac{4{,}8}{120}$; далее выполнены преобразования дроби и получен ответ $\frac1{25}$.
\block{Где именно возникает ошибка}
Первый неверный шаг --- составление дроби $\frac{4{,}8}{120}$. Во-первых, 4,8 записано в килограммах, а 120 --- в граммах: непосредственно делить эти числовые значения нельзя. Во-вторых, спрашивается масса \emph{котёнка к массе кошки}, значит, первым числом должно быть 120 г.
\block{Почему это неверно}
\errorword\ Отношение двух масс не зависит от единицы измерения, \emph{если обе массы выражены в одной и той же единице}. Дробь из чисел, подписанных разными единицами, не даёт искомого безразмерного отношения. Кроме того, перестановка числителя и знаменателя заменяет отношение на обратное.
\block{Ключевая идея}
\idea\ Для отношения именованных величин: 1) приведите их к одной единице; 2) сохраните указанный в вопросе порядок; 3) сократите дробь.
\block{Исправление и полное правильное решение}
Переведём массу кошки в граммы:
\[
4{,}8\text{ кг}=4{,}8\cdot1000\text{ г}=4800\text{ г}.
\]
Составим отношение массы котёнка к массе кошки:
\[
120:4800=\frac{120}{4800}=\frac{1}{40}.
\]
Здесь обе массы в граммах; разделили числитель и знаменатель на 120.
\block{Проверка}
$120\cdot40=4800$ г. Следовательно, кошка в 40 раз тяжелее котёнка, а масса котёнка составляет одну сороковую её массы.
\block{Итог}
\outcome\ Искомое отношение $\boxed{1:40}$, или $\boxed{\frac1{40}}$.
\block{Задача на закрепление}
Масса щенка --- 250 г, масса собаки --- 5 кг. Найдите отношение массы щенка к массе собаки.

\problem{7}{Отношение и прямая пропорция}
\block{Текст задания}
Для блюда смешивают 0,8 кг моркови и 5,6 кг капусты. Найдите: а) отношение моркови к капусте в виде $1:n$; б) массу капусты для 2,4 кг моркови при той же пропорции.
\block{Ваше решение и ответ}
На фотографии читается вычисление $\frac{0{,}8}{5{,}6}=\frac17$, помеченное буквой «б». Ни ответ $1:n$ для пункта «а», ни вычисление массы для пункта «б» не доведены до завершения. Зачёркнутые фрагменты не учитываются.
\block{Где именно возникает недочёт}
Последний правильный шаг --- нахождение отношения $\frac17$. Следующий обязательный шаг для пункта «а» --- запись $1:7$; он пропущен. Для пункта «б» необходимо использовать найденное соотношение и новую массу моркови, но соответствующего вычисления нет.
\block{Почему это неверно}
\errorword\ Дробь $\frac17$ численно верна, однако условие просит форму $1:n$. Кроме того, отношение --- лишь промежуточный результат: требуется определить массу капусты при изменении количества моркови.
\block{Ключевая идея}
\idea\ Если на одну часть моркови приходится семь таких же частей капусты, то масса капусты в 7 раз больше массы моркови при любом сохранении рецепта.
\block{Исправление и полное правильное решение}
\textbf{а)} Составим отношение моркови к капусте:
\[
0{,}8:5{,}6=1:7.
\]
Оба числа разделены на 0,8. Получили нужный вид $\boxed{1:7}$.

\textbf{б)} На 1 кг моркови приходится 7 кг капусты. Для 2,4 кг моркови необходимо:
\[
2{,}4\cdot7=\boxed{16{,}8\text{ кг капусты}}.
\]
Можно рассуждать иначе: 2,4 кг в 3 раза больше 0,8 кг, значит, капусты нужно $5{,}6\cdot3=16{,}8$ кг.
\block{Проверка}
$2{,}4:16{,}8=1:7$ --- отношение сохранилось.
\block{Итог}
\outcome\ а) $\boxed{1:7}$; б) $\boxed{16{,}8\text{ кг}}$.
\block{Задача на закрепление}
Для напитка смешивают 0,5 л сиропа и 2 л воды. Запишите отношение сиропа к воде в виде $1:n$ и найдите объём воды для 1,5 л сиропа при той же пропорции.

\problem{8}{Сравнение долей успеха}
\block{Текст задания}
Лена успешно выполнила 14 попыток из 20, а Маша --- 18 из 24. У кого доля удачных попыток больше?
\block{Ваше решение и ответ}
В решении записано, что у каждой было «6 успешных» попыток. Затем сравнены $\frac6{20}=\frac3{10}$ и $\frac6{24}=\frac14$; сделан вывод, что доля Лены больше.
\block{Где именно возникает ошибка}
Первый неверный шаг --- запись «успешных 6»: величины $20-14=6$ и $24-18=6$ являются количеством \emph{неудачных} попыток. Дальнейшие дроби сравнивают долю неудач, а не успехов.
\block{Почему это неверно}
\errorword\ Числитель должен отражать указанное в вопросе событие. У Лены успешно выполнены 14 попыток, у Маши --- 18. Если использовать 6, вычисляется доля неуспехов. При этом меньшая доля неудач означает большую долю успехов, поэтому вывод даже по смыслу получается противоположным.
\block{Ключевая идея}
\idea\ Сравнивать нужно относительные частоты одного и того же события: \emph{успехов с успехами}, а не абсолютные количества и не успехи с неудачами.
\block{Исправление и полное правильное решение}
Доля успешных попыток Лены:
\[
\frac{14}{20}=\frac7{10}=0{,}70=70\%.
\]
Доля успешных попыток Маши:
\[
\frac{18}{24}=\frac34=0{,}75=75\%.
\]
Поскольку $75\%>70\%$, доля удачных попыток больше у Маши.
\block{Проверка}
Проверим через доли неудач: у Лены $\frac6{20}=30\%$, у Маши $\frac6{24}=25\%$. Действительно, $100\%-30\%=70\%$, а $100\%-25\%=75\%$.
\block{Итог}
\outcome\ У Маши доля успехов больше: $\boxed{75\%>70\%}$.
\block{Задача на закрепление}
Саша выполнил 16 успешных бросков из 25, а Оля --- 21 из 30. У кого доля успешных бросков больше?

\problem{9}{Доля внутри подгруппы}
\block{Текст задания}
В секции 36\% ребят занимаются футболом, 24\% --- баскетболом, остальные --- лёгкой атлетикой. Какой процент \emph{среди занимающихся игровыми видами спорта} составляют футболисты?
\block{Ваше решение и ответ}
На предоставленных фотографиях решение задания №9 отсутствует. Ответ не указан.
\block{Где именно возникает недочёт}
Отсутствует обязательный первый шаг: определить новую совокупность для расчёта --- только тех ребят, которые занимаются игровыми видами спорта (футболом и баскетболом).
\block{Почему простого ответа «36\%» недостаточно}
\errorword\ Число 36\% характеризует долю футболистов \emph{среди всей секции}. В задаче спрашивается доля \emph{среди игроков}, то есть знаменатель изменился. Легкоатлетов из новой совокупности нужно исключить.
\block{Ключевая идея}
\idea\ Слова «среди ...» задают знаменатель. Сначала выясните размер этой группы, а затем найдите долю нужной категории внутри неё.
\block{Полное правильное решение}
Игровые виды спорта --- футбол и баскетбол. Вместе ими занимаются
\[
36\%+24\%=60\%
\]
всех ребят. Из этих 60 процентных пунктов 36 приходятся на футбол. Поэтому искомая доля:
\[
\frac{36}{60}=\frac35=0{,}6=\boxed{60\%}.
\]
Здесь 60\% --- доля \emph{игроков} среди всей секции, а 60\% в ответе --- доля \emph{футболистов среди игроков}. Совпадение чисел случайно: это разные величины.
\block{Проверка}
Баскетболисты среди игроков составляют $\frac{24}{60}=40\%$. Сумма $60\%+40\%=100\%$.
\block{Итог}
\outcome\ Среди занимающихся игровыми видами спорта футболисты составляют $\boxed{60\%}$.
\block{Задача на закрепление}
В кружках 45\% ребят занимаются шахматами, 30\% --- шашками, остальные --- рисованием. Какой процент среди занимающихся настольными играми составляют шахматисты?

\problem{10}{Сравнение долей и кратности}
\block{Текст задания}
В группе А успешны 18 попыток из 24, в группе Б --- 20 из 25. Определите: а) в какой группе выше доля успехов; б) во сколько раз большая доля превышает меньшую.
\block{Ваше решение и ответ}
В записи указаны числа 24 и 25 как общие количества, но «успешными» названы 6 и 5 попыток. Сравниваются $\frac6{24}=\frac14$ и $\frac5{25}=\frac15$; получены ответы «А» и $\frac54$.
\block{Где именно возникает ошибка}
Последний корректный шаг --- определение общих количеств 24 и 25. Первый неправильный шаг --- подстановка чисел $24-18=6$ и $25-20=5$ вместо числа успешных попыток. Эти числа означают \emph{неудачи}. Поэтому все дальнейшие вычисления относятся не к требуемым долям.
\block{Почему это неверно}
\errorword\ Формула доли успехов --- \emph{число успехов, делённое на число всех попыток}. Сравнение $\frac14$ с $\frac15$ математически выполнено верно, но оно отвечает на вопрос о доле неудач. Значение $\frac54$ показывает, во сколько раз доля неудач А больше доли неудач Б; это не ответ на пункт «б».
\block{Ключевая идея}
\idea\ Сначала вычислите обе доли \emph{успехов}. Для ответа «во сколько раз» большую из них разделите на меньшую.
\block{Исправление и полное правильное решение}
\textbf{а)} В группе А доля успехов:
\[
\frac{18}{24}=\frac34=0{,}75=75\%.
\]
В группе Б доля успехов:
\[
\frac{20}{25}=\frac45=0{,}8=80\%.
\]
Так как $80\%>75\%$, доля успехов выше в группе Б.

\textbf{б)} Большая доля равна $\frac45$, меньшая --- $\frac34$. Их отношение:
\[
\frac45:\frac34=\frac45\cdot\frac43
=\frac{16}{15}=1\frac1{15}.
\]
Значит, большая доля превышает меньшую в $\frac{16}{15}$ раза (примерно в $1{,}067$ раза).
\block{Проверка}
$\frac34\cdot\frac{16}{15}=\frac45$. Значит, множитель найден правильно. Разница долей --- 5 процентных пунктов, но для «во сколько раз» используется \emph{деление}, а не вычитание.
\block{Итог}
\outcome\ а) $\boxed{\text{группа Б}}$; б) $\boxed{\frac{16}{15}}$ раза.
\block{Задача на закрепление}
В группе А удачны 21 попытка из 28, а в группе Б --- 18 из 20. Определите, где доля успехов выше и во сколько раз большая доля превышает меньшую.

\vfill
\begin{center}
\rule{0.58\linewidth}{0.3pt}\\[-1pt]
{\footnotesize Лёвин Артём Александрович эксклюзивно для Ксении Васильченко}
\end{center}
\end{document}